% !TEX program = xelatex
\documentclass[aspectratio=169,12pt]{beamer}

% XeLaTeX用日本語設定
\usepackage{xltxtra}
\usepackage{zxjatype}
\usepackage[ipa]{zxjafont}

% パッケージ
\usepackage{amsmath,amssymb}
\usepackage{booktabs}
\usepackage{graphicx}
\usepackage{tikz}
\usepackage{pgfplots}
\pgfplotsset{compat=1.18}

% テーマ設定
\usetheme{Madrid}
\usecolortheme{seahorse}
\setbeamertemplate{navigation symbols}{}
\setbeamertemplate{footline}[frame number]

% タイトル
\title{ニューラルネットワークによる\\複数不可分財配分メカニズムの学習}
\subtitle{耐戦略性と個人合理性の下での効率的配分}
\author{B班}
\date{\today}

\begin{document}

%==============================================================================
\begin{frame}
\titlepage
\end{frame}

%==============================================================================
\section{問題設定}
%==============================================================================

\begin{frame}{不可能性定理と動機}
\textbf{Sönmez (1999) の不可能性定理}

複数財の交換問題において、以下の3条件を同時に満たすメカニズムは\\
\textbf{essentially single-valued core} に限られる：

\begin{enumerate}
    \item \textbf{耐戦略性 (SP)}：嘘の申告による利得が存在しない
    \item \textbf{個人合理性 (IR)}：取引後の効用が初期保有以上
    \item \textbf{効率性 (Efficiency)}：パレート改善の余地がない
\end{enumerate}
→具体例：TTC
\vspace{0.3cm}
\\\textbf{本研究のアプローチ}
\begin{itemize}
    \item IR・SPを制約として維持しつつ、\textbf{効率性ロスを最小化}
    \item ニューラルネットワークによる\textbf{Second-Best} メカニズムの学習
    \item 社会厚生（Welfare）の最大化を目的関数とする
\end{itemize}
\end{frame}

%==============================================================================
\section{モデル設定}
%==============================================================================

\begin{frame}{基本設定}
\textbf{エージェントと財}
\begin{itemize}
    \item エージェント：$N = \{1, 2, \ldots, n\}$（$n \in \{2, 3\}$）
    \item 財：$M = \{1, 2, \ldots, m\}$（$m \in \{5, 8\}$）
    \item 配分空間：$n = 2$ の場合 $2^m$ 通り、$n = 3$ の場合 $3^m$ 通り
    \item Disposal なし（全ての財を配分）
\end{itemize}

\vspace{0.3cm}
\textbf{選好パラメータ}
\begin{itemize}
    \item 財 $g$ への評価値：$v_{i,g} \sim \text{Uniform}(0, 1)$（iid）
    \item 補完財シナジー：$\alpha_i \sim \text{Uniform}(0, 0.5)$ (iid)
\end{itemize}

\vspace{0.3cm}
\textbf{初期保有}
\begin{itemize}
    \item 各財をランダムにエージェントに配分
    \item $\omega_i \subseteq M$：エージェント $i$ の初期保有バンドル
\end{itemize}
\end{frame}

%------------------------------------------------------------------------------
\begin{frame}{効用関数と制約}
\textbf{効用関数}（補完財シナジーを含む）

エージェント $i$ がバンドル $S_i \subseteq M$ を受け取る効用：
\[
u_i(S_i) = \sum_{g \in S_i} v_{i,g} + (|S_i| - 1) \cdot \alpha_i
\]
\vspace{0.2cm}
\textbf{制約条件}
\begin{block}{個人合理性 (IR)}
$\mathbb{E}[u_i(f(\theta))] \geq u_i(\omega_i) \quad \forall i \in N$ \hfill (ex-ante IR)
\end{block}

\begin{block}{耐戦略性 (SP)：正直申告が最適}
\[
\mathbb{E}[u_i(f(\theta_i, \theta_{-i}))] \geq \mathbb{E}[u_i(f(\theta'_i, \theta_{-i}))] \quad \forall \theta'_i, \forall i
\]
\begin{center}
SP違反 $= \max_{\theta'_i} \{\mathbb{E}[u_i(f(\theta'_i, \theta_{-i}))] - \mathbb{E}[u_i(f(\theta_i, \theta_{-i}))]\}$
\end{center}
\end{block}
\end{frame}

%------------------------------------------------------------------------------
\begin{frame}{目的関数}
\textbf{社会厚生の最大化}
\[
\max_{f} \sum_{i \in N} \mathbb{E}[u_i(f(\theta))]
\]
\begin{center}
s.t. IR ＆ SP
\end{center}

\vspace{0.5cm}
\textbf{Augmented Lagrangian による最適化}
\[
\mathcal{L} = -\text{Welfare} + \lambda_{\text{IR}} \cdot \text{IR違反} + \lambda_{\text{SP}} \cdot \text{SP違反} + \frac{\rho}{2}(\text{IR違反}^2+\text{SP違反}^2)
\]

\begin{itemize}
    \item $\lambda_{\text{IR}}, \lambda_{\text{SP}}$：双対変数
    \\{\footnotesize ※ 双対変数の更新：$\lambda \leftarrow \max(0, \lambda + \rho \cdot \text{違反量})$}
    \item $\rho$：ペナルティ係数
    \\{\footnotesize ※制約未達のときのみ徐々に増加}
\end{itemize}


\end{frame}

%==============================================================================
\section{実験設定}
%==============================================================================

\begin{frame}{実験条件}
\begin{table}[h]
\centering
\begin{tabular}{lccc}
\toprule
設定 & 2人5財 & 2人8財 & 3人5財 \\
\midrule
配分空間 & $2^5 = 32$ & $2^8 = 256$ & $3^5 = 243$ \\
学習ステップ & 20,000--30,000 & 20,000 & 20,000 \\
バッチサイズ & 512 & 512 & 512 \\
\bottomrule
\end{tabular}
\end{table}

\vspace{0.5cm}
\textbf{比較対象}
\begin{itemize}
    \item \textbf{Status quo}：取引なし（初期保有のまま）
    \item \textbf{Oracle}：制約なし/IR制約ありでの理論上限
    \item \textbf{TTC}：バンドル交換に限定したTop Trading Cycles
    \item \textbf{Greedy Swap}：初期配分からパレート改善的な単一財移動・交換を繰り返し、局所最適解に到達するまで実行する反復アルゴリズム
\end{itemize}
\end{frame}

%==============================================================================
\section{実験結果}
%==============================================================================

\begin{frame}{結果：2人5財}
\begin{table}[h]
\centering
\small
\begin{tabular}{lccc}
\toprule
手法 & Welfare & IR違反 & SP違反 \\
\midrule
Status quo & 3.26 & 0.000 & 0.000 \\
Oracle (制約なし) & \textbf{4.22} & -- & -- \\
Oracle (IR制約) & 3.85 & -- & -- \\
\midrule
Learned (lottery) & 3.50 & 0.062 & 0.017 \\
Learned (確定) & 3.53 & 0.755 & 0.120 \\
Bilateral TTC & 3.35 & 0.000 & 0.000 \\
\textbf{Greedy Swap} & \textbf{3.74} & 0.000 & 0.028 \\
\bottomrule
\end{tabular}
\end{table}

\vspace{0.3cm}
\begin{itemize}
    \item Learned は Status quo から \textbf{+0.24} の厚生改善
    \item ただし Greedy Swap が最も高い Welfare を達成
    \item Learned は IR/SP 違反が小さいが、完全には解消できていない
\end{itemize}
\end{frame}

%------------------------------------------------------------------------------
\begin{frame}{結果：2人8財}
\begin{table}[h]
\centering
\small
\begin{tabular}{lccc}
\toprule
手法 & Welfare & IR違反 & SP違反 \\
\midrule
Status quo & 5.52 & 0.000 & 0.000 \\
Oracle (制約なし) & \textbf{6.99} & -- & -- \\
Oracle (IR制約) & 6.67 & -- & -- \\
\midrule
Learned (lottery) & 5.76 & 0.097 & 0.020 \\
Learned (確定) & 5.78 & 0.943 & 0.374 \\
Bilateral TTC & 5.65 & 0.000 & 0.000 \\
\textbf{Greedy Swap} & \textbf{6.43} & 0.000 & 0.070 \\
\bottomrule
\end{tabular}
\end{table}

\vspace{0.3cm}
\begin{itemize}
    \item 8財では配分空間が $2^8 = 256$ に拡大
    \item Greedy Swap と Oracle の Gap は 0.56（制約なし）
    \item Learned は Status quo から \textbf{+0.24} の改善だが、Greedy に劣る
\end{itemize}
\end{frame}

%------------------------------------------------------------------------------
\begin{frame}{結果：3人5財}
\begin{table}[h]
\centering
\small
\begin{tabular}{lccc}
\toprule
手法 & Welfare & IR違反 & SP違反 \\
\midrule
Status quo & 3.10 & 0.000 & 0.000 \\
Oracle (制約なし) & \textbf{4.55} & -- & -- \\
Oracle (IR制約) & 3.91 & -- & -- \\
\midrule
\textbf{Learned (lottery)} & \textbf{3.49} & 0.052 & 0.007 \\
Learned (確定) & 3.50 & 0.414 & 0.069 \\
TTC Bundle Exchange & 3.29 & 0.000 & 0.000 \\
Greedy Swap & 3.10 & 0.000 & 0.000 \\
\bottomrule
\end{tabular}
\end{table}

\vspace{0.3cm}
\begin{itemize}
    \item 3人の場合、配分空間は $3^5 = 243$ 通り
    \item \textbf{Greedy Swap が Status quo と同等}（局所解に到達しにくい）
    \item Learned が \textbf{+0.40} の厚生改善で最良
\end{itemize}
\end{frame}

%==============================================================================
\section{考察}
%==============================================================================

\begin{frame}{結果の解釈}
\textbf{2人の場合：Greedy Swap優位}
\begin{itemize}
    \item 単一財の移動・交換でパレート改善を繰り返すだけで高い Welfare を達成
    \item 探索空間が小さく（$2^m$ 通り）、局所最適でも十分良い解に到達
    \item 結果として Learned (lottery) より高い Welfare を記録
\end{itemize}

\vspace{0.5cm}
\textbf{3人の場合：Learned (lottery) 優位}
\begin{itemize}
    \item 探索空間が複雑化（$3^m$ 通り）し、Greedy Swap はパレート改善を見つけにくい
    \item Greedy Swap は初期配分から動けず Status quo と同等に
    \item Learned (lottery) は大域的な配分を学習し、+0.4の厚生改善を達成
\end{itemize}

\vspace{0.5cm}
\textbf{TTCの解釈}
\begin{itemize}
    \item バンドル交換のみに制限した TTC は IR・SP を厳密に満たす
    \item 部分的な財の移動ができないため、Learned (lottery) より低い Welfareに
\end{itemize}
\end{frame}

%------------------------------------------------------------------------------
\begin{frame}{課題と限界}
\textbf{現状の課題}
\begin{itemize}
    \item Learned メカニズムは IR・SP 違反を完全には解消できていない
    \item $\lambda$ の調整が難しく、ペナルティが強すぎると安全な配分に収束
    \item 配分多様性が低い（unique allocations が少ない）
\end{itemize}

\vspace{0.5cm}
\textbf{Greedy Swap に勝ちにくい理由（仮説）}
\begin{enumerate}
    \item 問題設定のシンプルさ（補完財のみ、2人が多い）
    \item ニューラルネットの最適化が不十分
    \item そもそも Greedy Swap が理論的に強い可能性
\end{enumerate}
\end{frame}

%==============================================================================
\section{今後の拡張}
%==============================================================================

\begin{frame}{拡張の方向性}
\textbf{問題設定の複雑化}
\begin{itemize}
    \item \textbf{代替財の導入}：$\alpha_i < 0$（Disposal を許容）
    \item \textbf{望ましくない財}：負の評価値を持つ財
    \item \textbf{規模の拡大}：9財以上、4人以上
    \item \textbf{分布の変更}：正規分布、対数正規分布など
\end{itemize}

\vspace{0.5cm}
\textbf{手法の改善}
\begin{itemize}
    \item IR/SP ペナルティ等ハイパーパラメータの適応的調整
\end{itemize}

\vspace{0.5cm}
\textbf{理論的検討}
\begin{itemize}
    \item Greedy Swap の理論的性質の解明
\end{itemize}
\end{frame}

%==============================================================================
\begin{frame}{まとめ}
\begin{itemize}
    \item IR・SP を制約とし、Welfare を最大化するニューラルメカニズムを実装した
    \item 2人5財・8財、3人5財で実験を実施
    \item 2人の場合は Greedy Swap が強力なベースライン
    \item \textbf{3人以上では Learned メカニズムが優位}
    \item 今後は問題の複雑化により、学習メカニズムの優位性を検証
\end{itemize}
\end{frame}

%==============================================================================
\end{document}
